\documentclass[a4paper, serif, 10pt]{moderncv}

%% moderncv
% style and color
\moderncvstyle{banking}
\moderncvcolor{black}

%% page layout/margins
\usepackage[top=2.5cm, bottom=2.5cm, left=1.5cm, right=1.5cm]{geometry}

\nopagenumbers{}

%% Babel config for Dutch language
% ref:
% - https://latex3.github.io/babel/guides/locale-dutch.html
\usepackage[dutch]{babel}

%% fontspec
\usepackage{fontspec}

\IfFontExistsTF{Linux Libertine}{
  \setmainfont[Ligatures={TeX, Common}, Numbers=OldStyle]{Linux Libertine}
}{}
\IfFontExistsTF{Linux Libertine O}{
  \setmainfont[Ligatures={TeX, Common}, Numbers=OldStyle]{Linux Libertine O}
}{}

%% CTeX
% CJK support, used to show CJK characters in the resume
%
% - fontset=none: disable builtin fontset but instead set the CJK font manually
% - heading=false: disable ctex heading
% - punct=kaiming: use kaiming punctuations styles for CJK
% - scheme=plain: use plain scheme, do not override \normalsize font size
% - space=auto: space settings for CJK characters
%
% ref:
% - http://ctan.mirrorcatalogs.com/language/chinese/ctex/ctex.pdf
\usepackage[UTF8, heading=false, punct=kaiming, scheme=plain, space=auto]{ctex}

\IfFontExistsTF{Noto Serif CJK SC}{
  \setCJKmainfont{Noto Serif CJK SC}
}{}
\IfFontExistsTF{Noto Sans CJK SC}{
  \setCJKsansfont{Noto Sans CJK SC}
}{}

%% line spacing
\usepackage{setspace}
\setstretch{1.125}

%% URL styling - use normal text instead of monospace
\urlstyle{same}

% Auto-underline all links
% Defer until \begin{document} so hyperref is already loaded
\AtBeginDocument{%
  \let\oldhref\href
  \renewcommand{\href}[2]{\oldhref{#1}{\underline{#2}}}%
}

%% Patch \httplink and \httpslink to handle full URLs
% The original moderncv macros always prepend http:// or https:// to URLs.
% This causes issues when the URL already has a protocol (e.g., https://example.com)
% resulting in malformed URLs like https://https://example.com.
% This patch checks if the URL already contains "://" and uses it directly if so.
% Note: We use \str_set:Nx (x-expansion) to expand macros like \@homepage first.
\ExplSyntaxOn
\str_new:N \g_yamlresume_url_str

\RenewDocumentCommand{\httpslink}{O{}m}{
  \str_set:Nx \g_yamlresume_url_str {#2}
  \str_if_in:NnTF \g_yamlresume_url_str {://}
    { \url{#2} }
    { \url{https://#2} }
}

\RenewDocumentCommand{\httplink}{O{}m}{
  \str_set:Nx \g_yamlresume_url_str {#2}
  \str_if_in:NnTF \g_yamlresume_url_str {://}
    { \url{#2} }
    { \url{http://#2} }
}
\ExplSyntaxOff

\name{Sophie Bakker}{}
\title{Data Scientist met passie voor voorspellende modellen}
\phone[mobile]{+31 6 12345678}
\email{sophie.bakker@datavak.nl}
\homepage{https://www.sophiebakker.nl}

\address{Herengracht 123, Amsterdam, Noord-Holland, Nederland, 1011 AB}{}{}

\extrainfo{{\small \faLinkedin}\ \href{https://www.linkedin.com/in/sophiebakker}{@sophiebakker} {} {} {} • {} {} {}
{\small \faGithub}\ \href{https://github.com/sophiebakker}{@sophiebakker}}

\begin{document}

\maketitle

\section{Basis Informatie}

\cvline{}{Data Scientist met ruim 6 jaar ervaring in het ontwikkelen en implementeren van
datagedreven oplossingen voor complexe vraagstukken.
%
\begin{itemize}
\item Ontwerpt en implementeert voorspellende modellen voor klantverloop, vraagvoorspelling en fraudedetectie.
\item Combineert sterke analytische vaardigheden met een pragmatische, bedrijfsgerichte aanpak.
\item Ervaren met Python, R, SQL en cloudplatformen zoals AWS en Google Cloud.
\item Begeleidt junior datascientists en draagt bij aan de datastrategie van de organisatie.
\end{itemize}}
\section{Opleidingen}

\cventry{sep 2017–jul 2019}
        {Master, Data Science, Score: 8.2}
        {Universiteit van Amsterdam}
        {\url{https://www.uva.nl}}
        {}
        {\begin{itemize}
\item Afstudeeronderzoek naar causale inferentie in observationele studies, beoordeeld met een 9.
\item Specialisatie in voorspellende modellen en datagedreven besluitvorming.
\end{itemize}
\textbf{Cursussen}: Machine Learning, Statistical Learning, Big Data Technologies, Deep Learning, Data Visualisatie, Ethiek en Databeleid}

\cventry{sep 2014–jul 2017}
        {Bachelor, Kunstmatige Intelligentie, Score: 7.6}
        {Universiteit Utrecht}
        {\url{https://www.uu.nl}}
        {}
        {\begin{itemize}
\item Brede basis in AI-technieken, statistiek en softwareontwikkeling.
\item Minor Ondernemerschap gevolgd aan de Universiteit Utrecht.
\end{itemize}
\textbf{Cursussen}: Logica en Redeneren, Algoritmiek, Datastructuren, Probabilistische Modellen, Mens-Machine Interactie, Cognitiewetenschap}
\section{Werkervaring}

\cventry{jan 2022–Heden}
        {Senior Data Scientist}
        {Nova Analytics B.V.}
        {\url{https://www.nova-analytics.nl}}
        {}
        {\begin{itemize}
\item Geeft leiding aan de ontwikkeling van schaalbare machine-learning-pipelines voor klanten in de financiële sector.
\item Verbeterde het churnmodel met 15\% nauwkeurigheid door feature-engineering en modeloptimalisatie.
\item Implementeerde MLOps-praktijken zoals CI/CD en modelmonitoring, wat de time-to-market van modellen met 40\% verkortte.
\item Adviseert het management over datastrategie en bouwt aan een sterk data science-team.
\end{itemize}
\textbf{Trefwoorden}: Python, Machine Learning, MLOps, Strategie, Teamlead}

\cventry{sep 2019–dec 2021}
        {Data Scientist}
        {Zicht Data BV}
        {\url{https://www.zichtdata.nl}}
        {}
        {\begin{itemize}
\item Analyseerde grote datasets om logistieke processen efficiënter te maken.
\item Ontwikkelde realtime dashboards met Python, Streamlit en Tableau voor het monitoren van KPI's.
\item Werkte nauw samen met productowners en domeinexperts om databehoeften te vertalen naar analyses.
\item Voerde A/B-testen en statistische analyses uit ter ondersteuning van productbeslissingen.
\end{itemize}
\textbf{Trefwoorden}: Python, SQL, Data-analyse, Tableau, A/B-testen}

\cventry{feb 2019–jul 2019}
        {Data Analyst (stage)}
        {Datazicht Labs}
        {\url{https://www.datazichtlabs.nl}}
        {}
        {\begin{itemize}
\item Ondersteunde het team bij het opschonen, samenvoegen en analyseren van data voor beleidsonderzoek.
\item Realiseerde interactieve visualisaties met R en R Shiny.
\item Droeg bij aan datakwaliteitsrapportages en verbeterde de metadata van datasets.
\end{itemize}
\textbf{Trefwoorden}: R, Data Cleaning, Visualisatie, Rapportage}
\section{Talen}

\cvline{Nederlands}{Moedertaal of tweetalige vaardigheid \hfill \textbf{Trefwoorden}: Moedertaal}
\cvline{Engels}{Volledige professionele vaardigheid \hfill \textbf{Trefwoorden}: Cambridge C2, Zakelijk Engels}
\cvline{Duits}{Beperkte werkvaardigheid \hfill \textbf{Trefwoorden}: Basisvaardigheden}
\section{Vaardigheden}

\cvline{Machine Learning}{Expert \hfill \textbf{Trefwoorden}: Python, scikit-learn, TensorFlow, PyTorch, XGBoost, LightGBM}
\cvline{Statistiek}{Geavanceerd \hfill \textbf{Trefwoorden}: R, Hypothesetoetsing, Tijdreeksen, A/B-testen, Causale Inferentie}
\cvline{Data Engineering}{Intermediair \hfill \textbf{Trefwoorden}: SQL, Apache Spark, Airflow, Docker, dbt, AWS}
\cvline{Data Visualisatie}{Geavanceerd \hfill \textbf{Trefwoorden}: Tableau, Power BI, matplotlib, seaborn, ggplot2, Plotly, D3.js}
\cvline{Datamanagement}{Intermediair \hfill \textbf{Trefwoorden}: Data governance, Datakwaliteit, Metadata, AVG/GDPR}
\section{Onderscheidingen}

\cventry{nov 2019}
        {Nederlandse Vereniging voor Data Science}
        {Master Scriptieprijs Data Science}
        {}
        {}
        {Bekroond voor de beste afstudeerscriptie op het gebied van causale inferentie met observationele data.}

\cventry{jun 2021}
        {KLM}
        {KLM Datathon - Best Model}
        {}
        {}
        {Team behaalde de eerste plaats met een voorspellingsmodel voor bagageafhandelingsvertragingen.}
\section{Certificaten}

\cventry{mei 2022}
        {Amazon Web Services}
        {AWS Certified Machine Learning - Specialty}
        {\url{https://aws.amazon.com/certification/}}
        {}
        {}

\cventry{nov 2020}
        {Google}
        {Google Professional Data Engineer}
        {\url{https://cloud.google.com/certification/data-engineer}}
        {}
        {}
\section{Publicaties}

\cventry{feb 2022}
        {Een vergelijking van gradient boosting technieken voor vraagvoorspelling in de retailsector}
        {Tijdschrift voor Toegepaste Data Science}
        {\url{https://www.tijdschriftdata.nl/artikel/vergelijking-gradient-boosting}}
        {}
        {Onderzoekt de prestaties van XGBoost, LightGBM en CatBoost op retaildata en geeft praktische aanbevelingen voor het toepassen van deze technieken.}
\section{Projecten}

\cventry{mrt 2020–jul 2020}
        {Het project combineert tijdreeksanalyse met deep learning-modellen (LSTM) en biedt een API voor het ontsluiten van de voorspellingen. De resultaten worden gevalideerd aan de hand van publieke data van het Centraal Bureau voor de Statistiek.
}
        {Energieprognose voor slimme wijken}
        {\url{https://github.com/sophiebakker/energy-forecast}}
        {}
        {Open-source project dat energieverbruik op buurtniveau voorspelt op basis van historische data en weersvoorspellingen.
\textbf{Trefwoorden}: Python, Tijdreeksen, LSTM, Energietransitie}

\cventry{feb 2019–jun 2019}
        {Implementeerde een realtime fraudedetectiepijplijn met Apache Kafka voor streaming en scikit-learn voor de detectiemodellen. Evalueerde het systeem op een gesimuleerde transactiedataset en behaalde een F1-score van 0,91.
}
        {Fraudedetectiesysteem voor betaalverkeer}
        {\url{https://github.com/sophiebakker/fraud-detection}}
        {}
        {Eindproject van de masteropleiding: realtime fraudedetectiesysteem voor transacties.
\textbf{Trefwoorden}: Fraudedetectie, Kafka, scikit-learn, Realtime}
\section{Interesses}

\cvline{Wielrennen}{Toertochten, Amstel Gold Race}
\cvline{Schaken}{Clubcompetitie, Strategie}
\cvline{Lezen}{Sciencefiction, Vakliteratuur}
\cvline{Koken}{Vegetarisch, Wereldkeuken}
\section{Vrijwilligerswerk}

\cventry{jan 2020–Heden}
        {Data Scientist (vrijwilliger)}
        {Data for Good Nederland}
        {\url{https://www.dataforgoodnl.nl}}
        {}
        {\begin{itemize}
\item Verzorgt data-analyses voor non-profitorganisaties, waaronder de optimalisatie van voedselbankdistributie.
\item Ondersteunt vrijwilligers met het uitvoeren van dataprojecten.
\end{itemize}}

\end{document}